\section{Design-Based Partial-Document Supervision}\label{sec:min-audit}

The local laws encode the framework's key assumption: that the oracle
judgment of a long document decomposes across the leaves and merges of
a recursive summarization tree. Concretely, C1 says each leaf summary
preserves its span's evidence, C2 says re-summarizing a stored summary
does not drift, and C3 says split-then-merge preserves the same
cross-span interactions as summarize-together.

Two consequences of the assumption matter for the rest of the paper.
The first is root preservation (Theorem~\ref{thm:one-pass}): the
oracle judgment of the root summary equals the oracle judgment of the
full document. The second is schedule invariance
(Corollary~\ref{cor:schedule}): the choice of merge tree and chunk
grid does not change that judgment, so a researcher can move the
supervision granularity without changing the estimand. An equivalent
operator-side formulation appears in
Appendix~\ref{app:min-proofs}: the realized summarizer sits inside the
local-law subspace; equivalently, its approximation error vanishes on
that subspace.

The local laws are testable, and the audit is the procedure that
tests them. A realized tree is a finite population of partial-document
units. Leaves test C1. Stored summaries test C2. Merge nodes test C3.
Optional post-processing rounds test re-summary drift. A sampling design
selects units with logged positive propensities. An accessible judge
or gold oracle scores the sampled checks. Inverse-propensity weighting
transports those local observations back to the full population.

The design-based logic is the same one DSL uses for surrogate-label
inference. If we only inspect convenient nodes, we do not know what
population of local failures we measured. If we sample tree units at
random and log their inclusion probabilities, Horvitz--Thompson or
H\'ajek estimates recover the population signal under the design.
C-TreePO's extension is that the sampled units are leaves and merges,
the substructures that determine whether the labeled object still
represents the original document.

\begin{figure}[t]
\centering
\includegraphics[width=\linewidth]{01_base_sampled_grid.pdf}
\caption{\textbf{Six node-sample sets on the same tree.} Each panel
highlights one possible sampled set of tree units. Highlighted nodes
contribute to the local-law distortion estimate, weighted by their logged
propensities. The root can be sampled, but it is not privileged in the
node-level estimand.}
\label{fig:min-audit-sample}
\end{figure}

The audit reports two different quantities. \emph{Violation rates} are
diagnostic: which local law fails, and how often? \emph{Distortion means} are
certificate inputs: how far, in oracle units, is the realized tree from the
projection target where C1, C2, and C3 hold? A deployment should report both.
A low violation rate can still be costly if the failing nodes have large
distortion, and a high violation rate can be less damaging if each failure is
small in the task metric.

The natural design is nested. First sample documents from the target corpus.
Then sample nodes within selected documents. If the audit compares candidate
summaries or pairwise preferences, sample those actions within the selected
node. The logged probability for a labeled unit is the product of the relevant
design probabilities, for example
$p_i=p_d\,p_{u\mid d}$ for a node audit or
$p_i=p_d\,p_{u\mid d}\,p_{a\mid u}$ for an action-level label. Adaptive or
bandit-style sampling is allowed, but it needs an exploration floor and
logged propensities; otherwise the audit is no longer a design-based
estimator of the intended population.

For a realized audit population $\mathcal U$ of size $N$, unit $i$ has
inclusion indicator $Z_i$, logged propensity $\pi_i>0$, and distortion
contribution $Y_i$. The basic Horvitz--Thompson mean is
\[
  \hat\mu_{\mathrm{HT}}
  =
  \frac{1}{N}\sum_{i\in\mathcal U}\frac{Z_i}{\pi_i}Y_i.
\]
This is the node-level object. It is distinct from a root-observed corpus
evaluation, where every sampled document already has a full-document target
and the root prediction is scored directly against it.

The certificate has four pieces:
\[
  |G|
  \le
  C_{\mathrm{meth}}\hat{\Delta}_R
  + B_{\mathrm{cal}} + B_{\mathrm{est}} + B_{\mathrm{clip}}.
\]
$\hat{\Delta}_R$ is observed local distortion after $R$ rounds.
$C_{\mathrm{meth}}$ converts oracle distortion into the downstream objective
unit. $B_{\mathrm{cal}}$ covers the gap between the accessible judge and the
trusted target oracle. $B_{\mathrm{est}}$ is sampling error from querying only
some nodes. $B_{\mathrm{clip}}$ records stabilization bias from clipped
weights. More local labels shrink the design-based estimation term; they do
not remove calibration error between the judge and the social-science target.

This is where granularity becomes a methodological choice. If a paragraph,
section, or merge span can be audited against the same oracle-preservation
property as the full document, the researcher can observe cheaper local units
and aggregate them without changing the target. If local laws fail, the audit
does not hide that fact; it estimates the resulting distortion and shows
which law failed.

\subsection{Projection and Estimation: Two Uses of the Local Laws}\label{sec:min-projection-estimation}

Each local law is a statement about the oracle $f^*$ itself, not an
external constraint imposed on the learned operator. Written out:
\begin{align}
\text{C1:}\quad & \fstar(g(b)) \sim \fstar(b) && \text{on realized leaves } b, \notag \\
\text{C2:}\quad & \fstar(g(s)) \sim \fstar(s) && \text{on stored summaries } s \in \operatorname{range}(g), \notag \\
\text{C3:}\quad & \fstar\!\big(g(g(u)\concat g(v))\big) \sim \fstar(u\concat v) && \text{on merge calls } (u, v), \notag
\end{align}
all up to the oracle equivalence $\sim$ from
Section~\ref{sec:min-framework}. The framework uses these in two ways:
training-time projection and audit-time estimation.

\paragraph{Training-time projection.} The learned operator $(f, g)$
can be fit by minimizing
\begin{equation}\label{eq:min-J-lambda}
J_\lambda(f, g) = (1 - \lambda)\,L_{\text{oracle}}(f, g)
\;+\; \lambda\,L_{\text{local}}(f, g),
\end{equation}
where $L_{\text{oracle}}$ is the expected loss between $f$'s root
prediction and the observed root label, and $L_{\text{local}}$ is the
sum of expected per-law residuals on the realized calls of the
deployed tree. Adding the $\lambda$-weighted local-law term projects
the learned operator toward the subspace of $f^*$-approximations that
satisfy C1/C2/C3.

The endpoints have direct readings. $\lambda = 0$ is pure oracle
fit and the operator may sit anywhere in the function space; the
local laws play no role at training time. $\lambda = 1$ is pure
projection: the operator is pushed into the local-law subspace, with
no pressure to match $f^*$ on observed labels. Intermediate $\lambda$
trades the two off. The choice is a design decision rather than a
structural constant identified by the data: a practitioner with
strong prior confidence in the laws (e.g., classical mergeable
sketches whose merge is algebraically exact) can run with large
$\lambda$; a practitioner whose laws are an empirical bet on the
operator (the LLM case) sets $\lambda$ from validation behavior and
lets the audit report the residual the choice produced.

\paragraph{Audit-time estimation.} $L_{\text{local}}$ is the same
object the audit estimates on the deployed tree. The certificate's
$\hat{\Delta}_R$ above is the IPW-weighted estimate of $L_{\text{local}}$
at the realized call frequencies, computed from sampled units rather
than swept across the full population. Training projection and audit
estimation are pointed at the same target (distance from the
local-law subspace) from opposite sides of deployment; the
downstream consequence of the realized residual is the gap bound in
Section~\ref{sec:min-guarantees}.

The local laws reference an oracle $f^*$. The audit estimates $f^*$ at
sampled units from gold-standard labels collected by random assignment,
and computes local-law residuals against those estimates. Two valid
sources of gold-standard labels both work, and the audit math is the
same in either case.

\paragraph{Pre-existing root-level labels.} Document-level expert
codings exist in many social-science settings (Manifesto-Project
quasi-sentence codings are one example). Random assignment is satisfied
by sample splitting: hold out a randomly-chosen subset of documents at
training time and evaluate the audited tree against that held-out set.

\paragraph{Online node-level labeling.} When new gold-standard labels
are being collected during deployment, the audit consumes them at the
level they are collected. Random assignment is satisfied by drawing
audited nodes with logged inclusion probabilities.

The two modes are equivalent per unit of supervision in the controlled
Markov benchmark (Appendix~\ref{app:min-markov},
Figure~\ref{fig:min-markov-leaf-size}): tree supervision (root plus
sampled leaf labels) matches or improves on flat root-only supervision
across the budget range, with the gain depending mostly on the total
number of labeled units rather than the root-vs-leaf split. The
practical payoff: a researcher who has only document-level expert
labels can still audit, by sample splitting; a researcher who can pay
for new labels gets more signal per unit of expert effort by spending
those labels at the leaf level.

Appendix~\ref{app:min-audit} gives the design map, logging schema,
clustered uncertainty calculation, and Lean/proof anchors for the
DSL/IPW certificate layer.
